\documentclass[letterpaper, 12pt]{hmcpset}
\usepackage{cat-th-preamble}

\name{Jayden Thadani}
\class{Basic Category Theory self-study}
\assignment{Exercises 0 --- Introduction}

\begin{document}

\begin{problem}[0.10]
	Let $S$ be a set. The indiscrete topological space $\mathcal I(S)$ is the space whose set of points is $S$ and whose only open subsets are $\emptyset$ and $S$ itself. Imitating example 0.5, find a universal property satisfied by the space $\mathcal I(S)$.
\end{problem}

\begin{solution}
	Let $i$ be the inclusion that sends a point in $S$ to itself in $\mathcal I(S)$.\\\\
	Given any topological space $X$ and a function $F : X \to S$, the following diagram commutes.
	\[
		\begin{tikzcd}
			X \arrow[r, "F"] \arrow[rd, dashrightarrow, "\exists ! \, f \in \mathcal C({X, \, \mathcal I(S)})"']
			& S \arrow[d, "i"] \\
			& \mathcal I(S)
		\end{tikzcd}
	\]\\
	This is because we only have two open sets in $\mathcal I(S)$ with
	\[
		f^\text{pre}(\emptyset) = \emptyset
	\]
	and thus trivially open in $X$, and
	\[
		f^\text{pre}(S) = X
	\]
	(since every point in $X$ is mapped to some point in $S$), and thus, also trivially open in $X$.\\\\
	$f$ is obviously a unique function since any other continuous function that agrees with $i \circ F$ everywhere would have to agree with $f$ everywhere too. \\\\
	Thus, $\mathcal I(S)$ satisfies the universal property that for any topological space $X$ and function $F : X \to S$ there exists a unique continuous map $f : X \to \mathcal I(S)$ such that $f \circ i = F$.
\end{solution}

\begin{problem}[0.11]
	Fix a group homomorphism $\theta : G \to H$. Find a universal property satisfied by the pair $(\ker \theta, i)$.
\end{problem}

\begin{solution}
	Given a group $K$ and a homomorphism, $\varphi : K \to G$ such that ${\theta \circ \varphi = \varepsilon \circ \varphi}$, the following diagram commutes, with all the arrows being group homomorphisms.
	\[
		\begin{tikzcd}
			\ker \theta \arrow[r, hook, "i"]
				& G \arrow[r, "\theta", shift left] \arrow[r, "\varepsilon"', shift right]
				& H \\
			\forall K  \arrow[u, dashrightarrow, "\exists ! \, \varphi'"] \arrow[ur, "\varphi"]
		\end{tikzcd}
	\]
	That is, the pair $(\ker \theta, i)$ satisfy the universal property that for all groups $K$ with a homomorphism, $\varphi : K \to G$ such that ${\theta \circ \varphi = \varepsilon \circ \varphi}$, there is a homomorphism $\varepsilon' : K \to \ker \theta$ such that $i \circ \varphi' = \varphi$.\\\\
	This is because the condition on $\varphi$ forces us to have $\varphi^\text{img}(K) = \ker \theta$.
\end{solution}

\begin{problem}[0.12]
	Verify the universal property shown in diagram 0.3.
\end{problem}

\begin{solution}
	For a topological space $X$, covered by two open subsets, $U$ and $V$ and another topological space $Y$, we claim that if the diagram without the dashed arrow commutes, so does the diagram with the dashed arrow.
	\[
		\begin{tikzcd}
			U \cap V \arrow[r, hook, "i"] \arrow[d, hook, "j"] & U \arrow[d, hook, "j'"] \arrow[ddr, bend left, "\forall f"] \\
			V \arrow[r, hook, "i'"] \arrow[drr, bend right, "\forall g"] & X \arrow[dr, dashrightarrow, "\exists ! \, h"] \\
						&&\forall Y
		\end{tikzcd}
	\]\\
	This holds because if the diagram without the dashed arrow commutes, then $f$ and $g$ agree on $U \cap V$. Thus, we can choose $h(x) = f(x)$ for $x \in U$ and $h(x) = g(x)$ for $x \in V$ and still have $h(x)$ well defined for $x \in U \cap V$. The gluing lemma tells us that this stays continuous too.
\end{solution}

\begin{problem}[0.13]
	Denote by $\mathbb Z[x]$ the polynomial ring over $\mathbb Z$ in one variable.
	\begin{enumerate}
		\item Prove that for all rings $R$ and all $r \in R$, there exists a unique ring homomorphism $\varphi : \mathbb Z \to R$ such that $\varphi(x) = r$.
		\item Let $A$ be a ring and $a \in A$. Suppose that for all rings $R$ and all $r \in R$, there exists a unique ring homomorphism $\varphi : A \to R$ such that $\varphi(a) = r$. Prove that there is a unique isomorphism $\i : \mathbb Z[x] \to A$ such that $i(x) = a$.
	\end{enumerate}
\end{problem}

\begin{solution}
	\begin{enumerate}
		\item For any polynomial $p = \sum_{i = 1}^n a_i x^i$ consider the ring homomorphism given by
			\[
				\varphi(p) = \sum_{i = 1}^n \varphi(a_i) r^i.
			\]
			This satisfies the definition of a ring homomorphism and the requirement that ${\varphi(x) = r}$.\\\\
			We can show that this is unique by noticing that for any homomorphism $\psi$ fulfilling the requirements, we must have
			\[
				\begin{split}
					\psi(p) & = \psi \left ( \sum_{i = 1}^n a_i x^i \right ) \\
						& = \sum_{i = 1}^n \psi(a_i) \psi(x^i) \\
						& = \sum_{i = 1}^n \psi(a_i) \psi(x)^i \\
						& = \sum_{i = 1}^n \psi(a_i) r^i \\
						& = \varphi(p).
				\end{split}
			\]
		\item Suppose $A$ is the ring described. Then we have unique ring homomorphisms ${\varphi : A \to \mathbb Z[x]}$ and $\psi : \mathbb Z[x] \to A$ with $\varphi(a) = x$ and $\psi(x) = a$.\\\\ 
			Thus, $\psi \circ \varphi: A \to A$ is a ring homomorphism with $\psi \circ \varphi(a) = a$ and $\varphi \circ \psi : \mathbb Z[x] \to \mathbb Z[x]$ is a ring homomorphism with $\varphi \circ \psi(x) = x$. However, we already have such a ring homomorphisms --- $\operatorname{id}_A$ and $\operatorname{id}_{\mathbb Z[x]}$ respectively. Thus by the uniqueness condition, we must have $\psi \circ \varphi = \operatorname{id}_A$ and $\varphi \circ \psi = \operatorname{id}_{\mathbb Z[x]}$. Finally, this gives us that $\varphi$ and $\psi$ are both isomorphisms with $\psi = i$
	\end{enumerate}
\end{solution}

\begin{problem}[0.14]
	Let $X$ and $Y$ be vector spaces.
	\begin{enumerate}
		\item For the purposes of this exercise only, a \textit{cone} is a triple $(V, f_1, f_2)$ containing a vector space $V$, and linear maps $f_1 \in \mathcal L(V, X)$, $f_2 \in \mathcal L(V, Y)$. Find a cone $(P, p_1, p_2)$ with the following property: for all cones $(V, f_1, f_2)$, there exists a unique linear map $f : V \to P$ such that $p_1 \circ f = f_1$ and $p_2 \circ f = f_2$.
		\item Prove that there is essentially only one cone with the property stated in (a). That is, prove that if $(P, p_1, p_2)$ and $(P', p_1', p_2')$ both have this property then there is an isomorphism $i \in \mathcal L(P, P')$ such that $p_1' \circ i = p_1$ and $p_2' \circ i = p_2$.
		\item For the purposes of this exercise only, a \textit{cocone} is a triple $(V, f_1, f_2)$ consisting of a vector space $V$ and linear maps $f_1 \in \mathcal L(X, V)$, $f_2 \in \mathcal L(Y, V)$. Find a cocone $(Q, q_1, q_2)$ with the following property: for all cocones $(V, f_1, f_2)$, there exists a unique linear map $f : Q \to V$ such that $f \circ q_1 = f_1$ and $f \circ q_2 = f_2$.
		\item Prove that there is essentially only one cocone with the property stated in (c), in a sense that you should make precise.
	\end{enumerate}
\end{problem}

\begin{solution}
	\begin{enumerate}
		\item
			We draw the following diagram
			\[
				\begin{tikzcd}
					\forall V \ar[rrd, bend left, "f_1"] \ar[ddr, bend right, "f_2"] \ar[rd, dashrightarrow , "\exists ! \, f"] \\
					 & P \ar[r, "p_1"] \ar[d, "p_2"] & X \\
					 & Y 
				\end{tikzcd}
			\]
			Clearly $P = X \times Y$ with $p_1(x, y) = x$ and $p_2(x, y) = y$ as projection functions fit this diagram. For $f_1, f_2$, there is the unique linear map $f \in \mathcal L(V, P)$ with $v \mapsto (f_1 v, f_2 v)$.
		\item
			By the property from (a) that both cones have, there must be a unique linear map $p \in \mathcal L(P, P')$ with $p_j' \circ p = p_j$ ($j \in \mathbb N_2$). However, we must also have a unique linear map $p' \in \mathcal L(P', P)$ with $p_j \circ p' = p_j'$ (again, $j \in \mathbb N_2$). Thus, $p \circ p'$ is a linear map operator on $P$ with $p_j \circ (p \circ p') = p_j$ and $p' \circ p$ is a linear operator on $P'$ with $p_j' \circ (p' \circ p) = p_j'$ (for all $j \in \mathbb N_2$). However, we already have $\operatorname{id}_P$ and $\operatorname{id}_{P'}$ satisfying those, and thus, by the uniqueness condition we must have $p \circ p' = \operatorname{id}_P$ and $p' \circ p = \operatorname{id}_{P'}$. \\\\
			Thus, $p$ and $p'$ with $i = p$ being the desired isomorphism.\\\\
			We can see this isomorphism clearly in that the commutative diagrams can't tell $P$ and $P'$ apart:
			\[
				\begin{tikzcd}
					 P' \ar[rrd, bend left, "p_1'"] \ar[ddr, bend right, "p_2'"] \ar[rd, dashrightarrow , "\exists ! \, p'"] \\
					 & P \ar[r, "p_1"] \ar[d, "p_2"] & X \\
					 & Y 
				\end{tikzcd}
				\begin{tikzcd}
					 P \ar[rrd, bend left, "p_1"] \ar[ddr, bend right, "p_2"] \ar[rd, dashrightarrow , "\exists ! \, p"] \\
					 & P' \ar[r, "p_1'"] \ar[d, "p_2'"] & X \\
					 & Y 
				\end{tikzcd}
			\]

		\item
			Just reverse the arrows from (a)!
		\item
			Again, just reverse the arrows!
	\end{enumerate}
\end{solution}
	
\end{document}

